Para la implementación del sistema \textit{EKIS}, se ha seleccionado un stack tecnológico centrado en el ecosistema Oracle, priorizando la integridad de datos y la rapidez de desarrollo mediante el uso de Oracle APEX y PL/SQL.

\textbf{Diseño y Modelado de Datos}

El fundamento de la aplicación reside en un esquema de base de datos relacional robusto. En la transición del modelo conceptual al lógico, se aplicó la técnica de \textbf{fusión de tablas}, combinando estructuralmente múltiples tablas conceptuales para optimizar el rendimiento: las relacionadas con la gestión publicitaria se unificaron en la tabla \texttt{PUBLICACION}, y las de agrupación de comunicaciones en la tabla \texttt{MENSAJE}. El esquema final integra tablas clave como \textit{Usuarios}, \textit{Anuncios}, \textit{Publicaciones} y \textit{Hashtags}. Para garantizar la calidad de los datos, se configuraron restricciones de clave foránea con borrado en cascada y se implementaron validaciones tipo \texttt{CHECK} (valores 'Y'/'N') para suplir la ausencia de un tipo de dato booleano nativo.

\textbf{Lógica y Desarrollo en APEX}

La lógica del servidor se construyó utilizando el lenguaje PL/SQL, empleando procedimientos almacenados para operaciones transaccionales complejas y estandarizando la nomenclatura de variables (prefijo \texttt{p\_}) para asegurar el correcto funcionamiento. El desarrollo en Oracle APEX presentó desafíos técnicos específicos que definieron la arquitectura final:
\begin{itemize}
    \item \textbf{Procesamiento de Datos:} Inicialmente se intentó el uso de llamadas AJAX vía JavaScript, pero debido a fallos en la ejecución, se migró hacia procesos nativos del servidor (\textit{Processes}). Se configuró la ejecución de código SQL directo en lugar de \textit{Form Automatic Row Processing}, solucionando conflictos de persistencia y errores al modificar claves primarias.
    \item \textbf{Gestión de Memoria:} Se utilizaron cursores para iterar conjuntos de datos y \textbf{Colecciones APEX} (\texttt{APEX\_COLLECTION}) para la gestión temporal en memoria, vital para el algoritmo de listado de tendencias. Esto requirió una inicialización explícita de las colecciones para evitar errores en tiempo de ejecución.
    \item \textbf{Interfaz:} Se empleó el \textit{Universal Theme} para garantizar la adaptabilidad y se usaron componentes como \textit{List of Values} mapeadas para la asignación de anuncios.
\end{itemize}

\textbf{Seguridad y Control de Acceso}

La seguridad se gestionó mediante una arquitectura de roles (Usuario y Administrador) implementada a través de esquemas de autenticación y autorización personalizados:
\begin{description}
    \item[Autenticación:] Se desarrolló el paquete \texttt{PKG\_AUTH\_EKIS} para validar credenciales contra la tabla \texttt{USUARIO}, integrándolo en los componentes compartidos de la aplicación.
    \item[Autorización:] Se restringió el acceso a funcionalidades críticas, como la asignación de anuncios, mediante lógica condicional en la interfaz y redirecciones de URL personalizadas. Esto asegura que los requisitos funcionales administrativos permanezcan ocultos para los usuarios estándar.
\end{description}